% GENERATED FILE. Edit canonical lessons, then run npm run generate:books.
% canonical-source-hash: fnv1a64:9509b7b0ae36f865
% canonical-lessons: RU-C475-vostorg, RU-C475-voskhishchatsya, RU-C475-gorditsya, RU-C475-gordost, RU-C475-khvastatsya

\chapter{Delight and Pride}
\label{ch:ru-delight-and-pride}
\hlchaptermodality{475}

\begin{chapteropening}
\textbf{By the end of this chapter:} \emph{I can talk about delight, admiration and pride.}

Complete the last lesson of chapter 475: I can talk about delight, admiration and pride.
\end{chapteropening}

\section[\texorpdfstring{\ru{восторг} (vostórg)}{восторг (vostórg)}]{\ru{восторг} (vostórg) — delight, rapture}

\label{lesson:RU-C475-vostorg}



\begin{quote}
Before the new one: say the Russian for worry, then the Russian for to calm down.
\end{quote}

\subsection*{You'll want to know: \ru{восторг}}
\textbf{\ru{восторг}} (\emph{vostórg}) — \textquotedblleft{}delight, rapture\textquotedblright{}.

\begin{quote}
\ru{Дети} \ru{были} \ru{в} \ru{восторге} \ru{от} \ru{цирка}.
\end{quote}

\begin{quote}
\emph{Déti býli v vostórge ot tsírka.}
\end{quote}

\begin{quote}
The children were delighted with the circus.
\end{quote}

Masculine. \emph{\ru{Быть} \ru{в} \ru{восторге} \ru{от}} + genitive — to be thrilled with.

\begin{cousinweb}
\emph{\ru{восторг}} is built on \emph{\ru{вос}-} 'up' and a root usually linked to the one in \emph{\ru{исторгнуть}} 'to tear out', so rapture is being carried away; EtymDB has no entry, so the link is not settled.
\end{cousinweb}

\begin{grammarlens}[title={What you've built}]
One more word for telling, explaining and describing at B1.
\end{grammarlens}

\subsection*{Your turn}
\begin{itemize}
\raggedright
  \item \emph{Say it:} \emph{vostórg}
  \item \emph{Say it:} \emph{vostórg}, once more
  \item \emph{Say it:} a sentence of your own with \emph{vostórg}
  \item \emph{Recall:} say the Russian for worry, then the Russian for to calm down, then say \emph{vostórg} again
\end{itemize}

\begin{tcolorbox}[breakable,colback=teal!4,colframe=teal!35!black,title={Before you move on}]
What does \ru{восторг} mean? (\textquotedblleft{}Delight, rapture\textquotedblright{}.) Say it once more.
\end{tcolorbox}
\section[\texorpdfstring{\ru{восхищаться} (voskhishchátsya)}{восхищаться (voskhishchátsya)}]{\ru{восхищаться} (voskhishchátsya) — to admire}

\label{lesson:RU-C475-voskhishchatsya}



\begin{quote}
Before the new one: say the Russian for to calm down, then the Russian for delight.
\end{quote}

\subsection*{You'll want to know: \ru{восхищаться}}
\textbf{\ru{восхищаться}} (\emph{voskhishchátsya}) — \textquotedblleft{}to admire\textquotedblright{}.

\begin{quote}
\ru{Я} \ru{восхищаюсь} \ru{её} \ru{терпением}.
\end{quote}

\begin{quote}
\emph{Ya voskhishcháyus yeyó terpéniyem.}
\end{quote}

\begin{quote}
I admire her patience.
\end{quote}

Imperfective; perfective \ru{восхититься} (\ru{восхищусь}). Takes the instrumental: \emph{\ru{восхищаться} \ru{кем}, \ru{чем}}.

\begin{cousinweb}
\emph{\ru{восхищаться}} is built on \emph{\ru{восхитить}}, which was borrowed from Church Slavonic \emph{v\ru{ъ}sxytiti} 'to snatch away, to fascinate' (EtymDB): admiration carries you away.
\end{cousinweb}

\begin{grammarlens}[title={What you've built}]
One more word for telling, explaining and describing at B1.
\end{grammarlens}

\subsection*{Your turn}
\begin{itemize}
\raggedright
  \item \emph{Say it:} \emph{voskhishchátsya}
  \item \emph{Say it:} \emph{voskhishchátsya}, once more
  \item \emph{Say it:} a sentence of your own with \emph{voskhishchátsya}
  \item \emph{Recall:} say the Russian for to calm down, then the Russian for delight, then say \emph{voskhishchátsya} again
\end{itemize}

\begin{tcolorbox}[breakable,colback=teal!4,colframe=teal!35!black,title={Before you move on}]
What does \ru{восхищаться} mean? (\textquotedblleft{}To admire\textquotedblright{}.) Say it once more.
\end{tcolorbox}
\section[\texorpdfstring{\ru{гордиться} (gordítsya)}{гордиться (gordítsya)}]{\ru{гордиться} (gordítsya) — to be proud}

\label{lesson:RU-C475-gorditsya}



\begin{quote}
Before the new one: say the Russian for delight, then the Russian for to admire.
\end{quote}

\subsection*{You'll want to know: \ru{гордиться}}
\textbf{\ru{гордиться}} (\emph{gordítsya}) — \textquotedblleft{}to be proud\textquotedblright{}.

\begin{quote}
\ru{Отец} \ru{гордится} \ru{своим} \ru{сыном}.
\end{quote}

\begin{quote}
\emph{Otéts gordítsya svoím sýnom.}
\end{quote}

\begin{quote}
The father is proud of his son.
\end{quote}

Imperfective; present \ru{горжусь}, \ru{гордишься} (\ru{д} becomes \ru{ж}). Takes the instrumental: \emph{\ru{гордиться} \ru{сыном}}.

\begin{cousinweb}
\emph{\ru{гордиться}} is built on \emph{\ru{гордый}} 'proud', which is inherited from Proto-Slavic \emph{g\ru{ъ}rd\ru{ъ}} 'proud' (EtymDB).
\end{cousinweb}

\begin{grammarlens}[title={What you've built}]
One more word for telling, explaining and describing at B1.
\end{grammarlens}

\subsection*{Your turn}
\begin{itemize}
\raggedright
  \item \emph{Say it:} \emph{gordítsya}
  \item \emph{Say it:} \emph{gordítsya}, once more
  \item \emph{Say it:} a sentence of your own with \emph{gordítsya}
  \item \emph{Recall:} say the Russian for delight, then the Russian for to admire, then say \emph{gordítsya} again
\end{itemize}

\begin{tcolorbox}[breakable,colback=teal!4,colframe=teal!35!black,title={Before you move on}]
What does \ru{гордиться} mean? (\textquotedblleft{}To be proud\textquotedblright{}.) Say it once more.
\end{tcolorbox}
\section[\texorpdfstring{\ru{гордость} (górdost)}{гордость (górdost)}]{\ru{гордость} (górdost) — pride}

\label{lesson:RU-C475-gordost}



\begin{quote}
Before the new one: say the Russian for to admire, then the Russian for to be proud.
\end{quote}

\subsection*{You'll want to know: \ru{гордость}}
\textbf{\ru{гордость}} (\emph{górdost}) — \textquotedblleft{}pride\textquotedblright{}.

\begin{quote}
\ru{Она} \ru{с} \ru{гордостью} \ru{показала} \ru{нам} \ru{свой} \ru{диплом}.
\end{quote}

\begin{quote}
\emph{Oná s górdostyu pokazála nam svoy diplóm.}
\end{quote}

\begin{quote}
She proudly showed us her diploma.
\end{quote}

Feminine. \emph{\ru{С} \ru{гордостью}} — proudly.

\begin{cousinweb}
\emph{\ru{гордость}} is built on \emph{\ru{гордый}} 'proud' with \emph{-\ru{ость}}; \emph{\ru{гордый}} is inherited from Proto-Slavic \emph{g\ru{ъ}rd\ru{ъ}} (EtymDB).
\end{cousinweb}

\begin{grammarlens}[title={What you've built}]
One more word for telling, explaining and describing at B1.
\end{grammarlens}

\subsection*{Your turn}
\begin{itemize}
\raggedright
  \item \emph{Say it:} \emph{górdost}
  \item \emph{Say it:} \emph{górdost}, once more
  \item \emph{Say it:} a sentence of your own with \emph{górdost}
  \item \emph{Recall:} say the Russian for to admire, then the Russian for to be proud, then say \emph{górdost} again
\end{itemize}

\begin{tcolorbox}[breakable,colback=teal!4,colframe=teal!35!black,title={Before you move on}]
What does \ru{гордость} mean? (\textquotedblleft{}Pride\textquotedblright{}.) Say it once more.
\end{tcolorbox}
\section[\texorpdfstring{\ru{хвастаться} (khvástatsya)}{хвастаться (khvástatsya)}]{\ru{хвастаться} (khvástatsya) — to boast, to show off}

\label{lesson:RU-C475-khvastatsya}



\begin{quote}
Before the new one: say the Russian for to be proud, then the Russian for pride.
\end{quote}

\subsection*{You'll want to know: \ru{хвастаться}}
\textbf{\ru{хвастаться}} (\emph{khvástatsya}) — \textquotedblleft{}to boast, to show off\textquotedblright{}.

\begin{quote}
\ru{Он} \ru{любит} \ru{хвастаться} \ru{своей} \ru{машиной}.
\end{quote}

\begin{quote}
\emph{On lyúbit khvástatsya svoyéy mashínoy.}
\end{quote}

\begin{quote}
He likes to show off his car.
\end{quote}

Imperfective; the perfective is \ru{похвастаться}. Takes the instrumental: \emph{\ru{хвастаться} \ru{машиной}}.

\begin{cousinweb}
The origin of \emph{\ru{хвастаться}} is not settled: EtymDB lists the word without a source, and it is often compared with \emph{\ru{хвалить}} 'to praise'.
\end{cousinweb}

\begin{grammarlens}[title={What you've built}]
One more word for telling, explaining and describing at B1.
\end{grammarlens}

\subsection*{Your turn}
\begin{itemize}
\raggedright
  \item \emph{Say it:} \emph{khvástatsya}
  \item \emph{Say it:} \emph{khvástatsya}, once more
  \item \emph{Say it:} a sentence of your own with \emph{khvástatsya}
  \item \emph{Recall:} say the Russian for to be proud, then the Russian for pride, then say \emph{khvástatsya} again
\end{itemize}

\begin{tcolorbox}[breakable,colback=teal!4,colframe=teal!35!black,title={Before you move on}]
What does \ru{хвастаться} mean? (\textquotedblleft{}To boast, to show off\textquotedblright{}.) Say it once more.
\end{tcolorbox}
